\section{Introduction}
\label{sec:introduction}

Consider a clinical decision support system that alerts physicians to potential myocardial
infarction. The system reports a risk probability of 0.80 for every patient who eventually
receives an MI diagnosis --- but it also reports 0.80 for 30\% of healthy patients. Its AUROC
is excellent; it correctly ranks patients by risk. Yet its probabilities are useless for
threshold-based decision making, because a physician who trusts the 0.80 output will either
over-treat or under-treat depending on the threshold. AUROC is blind to this failure.
Probability calibration is not. This gap --- between how models are benchmarked and how they
are used clinically --- is the problem this paper addresses.

ECG foundation models (FMs) have rapidly expanded the capability of automated cardiac
analysis. Models such as ST-MEM~\citep{na2024stmem}, MERL~\citep{liu2024merl},
HuBERT-ECG~\citep{hubertecg2024}, ECGFM-KED~\citep{ecgfmked2024}, and
ECG-FM~\citep{mckeen2024ecgfm} leverage large-scale ECG pretraining to produce frozen
representations that transfer to downstream classification tasks via simple linear probes.
The BenchECG evaluation suite~\citep{lunelli2025benchecg} and the ``beyond accuracy''
benchmark~\citep{filice2026beyond} have standardized AUROC and F1 comparisons across
these models. Yet neither benchmark systematically reports calibration metrics (ECE, Brier
score, signed calibration error) or post-hoc recalibration results. When models are
deployed for multi-label cardiac risk stratification --- where probability outputs feed
directly into alert thresholds --- this omission is a patient-safety concern.

\paragraph{What we found.}
This paper presents the first calibration audit of ECG FMs, spanning 24 model$\times$dataset
cells across four public datasets. Four findings emerged from the data:

\begin{enumerate}
  \item \textbf{Higher discrimination, worse calibration on PTB-XL.} ST-MEM achieves the
    highest AUROC (0.838) yet poor calibration (ECE 0.124), and the masked-prediction
    HuBERT-ECG is the worst calibrated (ECE 0.136). All five foundation models beat the
    supervised S4D baseline on AUROC but trail its calibration (ECE 0.064) by 0.015--0.072.
    The discrimination--calibration tradeoff is systematic, not random.

  \item \textbf{The tradeoff disappears on single-label arrhythmia tasks.} On CPSC2018,
    CODE-15\%, and MIMIC-IV-ECG, all six encoders are comparatively well-calibrated and
    high-AUROC models are \emph{also} well-calibrated. Task structure---multi-label vs.\
    single-label---rather than model architecture drives calibration difficulty.

  \item \textbf{Within-domain pretraining is not sufficient.} On MIMIC-IV-ECG, ECG-FM---
    which was \emph{pretrained on MIMIC-IV-ECG}---reaches only AUROC 0.771, behind MERL
    (0.861) and ST-MEM (0.854) in all 10 resplits, and merely on par with the supervised
    S4D baseline (0.793; the two are statistically tied). Backbone domain alignment does not
    guarantee the best linear-probe performance.

  \item \textbf{Pretraining objective leaves a calibration signature.} MERL (contrastive
    InfoNCE) shows the strongest underconfidence on PTB-XL ($\text{SCE}=-0.003$);
    masked-prediction models (ST-MEM, HuBERT-ECG) are near-neutral ($\text{SCE}\approx0$);
    and knowledge-enhanced, hybrid, and supervised models are mildly overconfident
    ($\text{SCE}=+0.007$--$+0.008$). Feature geometry tracks this: MERL's contrastive
    objective yields extreme feature spread (std of L2 norms 34.8 vs.\ 0.23 for ST-MEM),
    while HuBERT-ECG has the highest class separation (0.046) and the highest ECE.
\end{enumerate}

\paragraph{Contributions.}
\begin{enumerate}
  \item \textbf{Calibration benchmark:} A 24-cell AUROC--ECE--SCE--Brier evaluation of five
    ECG FMs and a supervised S4D baseline across PTB-XL, CPSC2018, CODE-15\%, and
    MIMIC-IV-ECG, with patient-level splits, per-class multi-label calibration evaluation,
    minimum-support filtering, and bootstrap 95\% confidence intervals on every headline
    metric. All results are machine-readable and reproducible from cached features.

  \item \textbf{Pareto frontier and composite scoring:} We formally define the AUROC--ECE
    Pareto frontier and a composite score $C(\alpha) = \alpha \cdot \text{AUROC}
    + (1-\alpha)(1 - \text{ECE}/\text{ECE}_{\max})$ that lets practitioners trade off
    discrimination and calibration explicitly. On PTB-XL, \textbf{S4D ranks first} at
    $\alpha=0.7$ (score 0.699) and MERL at $\alpha=0.9$ (0.781), while ST-MEM---the AUROC
    leader---falls to 5th of 6 (0.614) and HuBERT-ECG ranks \emph{last} (0.551). After
    isotonic recalibration the ordering inverts: ST-MEM rises to \emph{first} ($C(0.7)=0.789$),
    showing that recalibration changes the recommended model. On CODE-15\%, CPSC2018, and
    MIMIC-IV-ECG, MERL, ST-MEM, and S4D form the Pareto front.

  \item \textbf{RankSafe-Cal:} A recalibration wrapper that selects, per class, the
    lowest-ECE calibrator whose out-of-sample AUROC loss is within a tolerance $\epsilon$,
    turning ``calibration can hurt AUROC'' into an explicit constraint. It recovers most of
    isotonic's ECE reduction (PTB-XL $0.097\!\to\!0.054$) at $\sim$6$\times$ smaller AUROC cost,
    and on MIMIC-IV-ECG dominates unconstrained isotonic on both axes
    (Section~\ref{sec:ranksafe}).

  \item \textbf{Reliability stress tests:} (i) Cross-dataset calibration \emph{transfer}---a
    calibrator ported across cohorts succeeds in only 9--42\% of endpoint/model pairs, so
    calibration is largely deployment-context-specific; and (ii) split-conformal prediction,
    which attains its $90\%$ coverage guarantee on every endpoint and complements recalibration
    with finite-sample coverage. We also add proper scores (NLL), calibration slope/intercept,
    a mechanistic ECE regression, multi-seed robustness, and all-model subgroup audits.

  \item \textbf{PhysioCalib:} A severity-weighted isotonic calibration method that allocates
    more correction to high-stakes conditions (MI, STTC) and less to routine findings (NORM),
    formally a monotone rank-preserving method, implemented in the \texttt{ecg-audit} package.

  \item \textbf{ecg-audit:} A pip-installable, unit-tested Python package (open-sourced at
    \url{https://github.com/amensk/ecg-audit}) providing one-command calibration auditing,
    the recalibration methods above, decision-curve analysis, subgroup calibration,
    reliability diagrams, and a community leaderboard scaffold.
\end{enumerate}

\paragraph{Paper organisation.}
Section~\ref{sec:related_work} reviews ECG FM benchmarks and calibration literature.
Section~\ref{sec:method} defines all metrics, the Pareto frontier, composite score,
PhysioCalib, and the decision-curve analysis framework.
Section~\ref{sec:experiments} describes datasets, models, and splits.
Section~\ref{sec:results} presents the main benchmark results.
Section~\ref{sec:clinical_utility} presents decision-curve analysis and PhysioCalib results.
Section~\ref{sec:open_source} describes the \texttt{ecg-audit} package.
Section~\ref{sec:discussion} discusses mechanisms and implications.
Section~\ref{sec:conclusion} concludes.
